\chapter{The burning cell}
\label{ch:burning-cell}

\checkedagainst{1.4.1 (\code{7092d89})}{The closure's closed forms are
derived symbolically in \code{doc/manual/checks/ch6_closure.py}; the code
is checked against 50-digit values by \code{test_hvrb}.}

A burn model tells a cell how fast to react. Everything else a reacting
cell needs (what its pressure, sound speed and temperature are, part way
through the reaction) is the same for every model in \bw, and is the
subject of this chapter. \Cref{ch:burn-models} then gives each model's rate.

\section{What a burn adds to a cell}
\label{sec:burn-adds}

A material that burns (deck: \key{burn = hvrb} or \key{thermalcascade},
with \key{eos = mie_gruneisen_t} and \key{products = ideal_gas}) carries one
more number in every cell, the \emph{extent of reaction} $\lam$: the mass
fraction of the cell that has become detonation products, from $\lam = 0$
(unreacted) to $\lam = 1$ (burnt out). On a Lagrangian mesh the cell
carries its own $\lam$ with it, so there is nothing to advect. Three things
are added to the hydrodynamics, and nothing else in the cycle changes:
\begin{enumerate}
\item a \emph{rate law}, which moves $\lam$ forward in time (each model's
  own, \cref{ch:burn-models});
\item a \emph{closure}, which gives the cell's pressure and sound speed from
  its density $\rho$, specific internal energy $e$ and $\lam$
  (\cref{sec:closure,sec:sound-speed});
\item a place in the cycle where $\lam$ is updated (\cref{sec:one-cycle}).
\end{enumerate}

The chemical energy is carried in $e$ itself: the products' energy is
measured so that reaction releases the heat of reaction $Q$ into the
mixture's temperature, at fixed $\rho$ and $e$ (\cref{sec:products}).
Burning therefore does no work and adds no energy of its own: total energy
is conserved exactly as in an inert run, and the run's energy error (the
heartbeat file) is the check.

\section{The unreacted solid}
\label{sec:solid}

The reactant is a Mie--Gr\"uneisen solid with a temperature (deck:
\key{eos = mie_gruneisen_t}, with the specific heat \key{CV} and the
temperature at rest \key{T0}). With $\rho\Gamma = \rhoz\Gz$ its pressure is
\begin{equation}
  \Ps(\rho, e) = \rhoz\Gz\, e + \Pi(\rho),
  \label{eq:mg}
\end{equation}
where $\Pi$ depends on the density alone: the Hugoniot's part in
compression, and a cubic in $\rho - \rhoz$ in expansion (\cref{ch:materials}).
\trace{modules/eos_mod.f95}{mg_pressure}{test_eos}

The temperature needs two curves of the volume alone. With the
compression $\eta = 1 - \rhoz/\rho$, the \emph{isentrope from rest},
$e_S(\rho)$, is the solution of
\begin{equation}
  \frac{\dd e_S}{\dd\rho} = \frac{\Ps(\rho, e_S)}{\rho^2},
  \qquad e_S(\rhoz) = 0,
  \label{eq:isentrope}
\end{equation}
and its temperature is $T_S = T_0\exp(\Gz\eta)$. A state off the isentrope
at the same density differs from it by $\Cv$ times the temperature
difference, $e = e_S + \Cv(T - T_S)$. It is convenient to collect the
temperature-independent part into the \emph{zero-temperature curve},
\begin{equation}
  e_c(\rho) = e_S(\rho) - \Cv T_S(\rho), \qquad
  \Pc(\rho) = \Pi(\rho) + \rhoz\Gz\, e_c(\rho),
  \label{eq:cold}
\end{equation}
the energy and pressure the solid would have at $T = 0$, so that, with
$\thv \equiv \Cv T$,
\begin{equation}
  e = e_c(\rho) + \thv, \qquad
  \Ps = \Pc(\rho) + \rhoz\Gz\,\thv, \qquad
  \frac{\dd e_c}{\dd\rho} = \frac{\Pc}{\rho^2}.
  \label{eq:solid-theta}
\end{equation}
\Cref{eq:solid-theta} is the form the closure uses: the solid's state is
its zero-temperature curve plus a thermal part $\thv$, linear in $T$.
\trace{modules/temp_mod.f95}{temp_cold_cell, cold_step, temp_of}{test_temp}

\paragraph{How a cell carries $e_S$.} \Cref{eq:isentrope} is linear in
$e_S$, so its solution from one compression to the next is exact, except for
one integral of $\Pi$. \bw\ evaluates that integral by three-point
Gauss--Legendre quadrature, on pieces no longer than $\Delta\eta = 0.005$. Each cell carries
its own $e_S$ and moves it to its new density once a step, rather than
reading a table, so that constants graded cell by cell (planned) need no
table per material. (Module \code{temperature} calls this carried value
\code{ec}; module \code{mixture} calls $e_c$ of \cref{eq:cold}
\code{ecs}. This manual uses $e_S$ and $e_c$.)

\section{The products}
\label{sec:products}

The detonation products are an ideal gas (deck: \key{products =
ideal_gas}, with \key{GAMMA} and \key{Q}) with the solid's specific heat:
\begin{equation}
  \Pg = (\gamma - 1)\rho\,\Cv T = (\gamma - 1)\rho\,\thv, \qquad
  e_{\mathrm{gas}} = \Cv T - Q, \qquad
  c_{\mathrm{gas}}^2 = \gamma(\gamma - 1)\thv.
  \label{eq:gas}
\end{equation}
At $T = 0$ the products' energy is $-Q$, the heat of reaction below the
solid's zero, so a cell that reacts at fixed $e$ gains $Q$ of thermal energy
per unit mass reacted. One $\Cv$ for both phases is a simplification (real products
have a specific heat of their own); it is what makes the closure below exact
without iteration. Another products equation of state (JWL, a
table) would be another set of routines in module \code{products}, and with
its own $\Cv$ the closure would become a one-variable solve for $T$ in the
same place.
\trace{modules/prod_mod.f95}{idg_state}{test_hvrb}

\section{The closure: one density, one temperature}
\label{sec:closure}

A partly burnt cell is a mixture of solid and gas. \bw\ puts both phases at
the cell's density and at one temperature (\emph{density--temperature
equilibrium}). The phases' pressures then differ, and the cell's pressure
is their mass-weighted mean. (Pressure--temperature equilibrium, the common
alternative, solves for two densities instead.) Because both phases have
the same density, the products' volume fraction equals their mass fraction:
$\lam$ is both.

The mixture's energy is the mass-weighted sum of the phases', by
\cref{eq:solid-theta,eq:gas}:
\begin{equation*}
  e = (1 - \lam)\,(e_c + \thv) + \lam\,(\thv - Q).
\end{equation*}
It is linear in $\thv$, so given $(\rho, e, \lam)$ the temperature follows
at once, and with it the pressure:
\begin{align}
  \thv &= e - (1 - \lam)\, e_c(\rho) + \lam Q, \qquad T = \thv/\Cv,
  \label{eq:theta}\\
  P &= (1 - \lam)\,\bigl[\Pc(\rho) + \rhoz\Gz\,\thv\bigr]
     + \lam\,(\gamma - 1)\rho\,\thv.
  \label{eq:pmix}
\end{align}
\trace{modules/mix_mod.f95}{mix_state}{test_hvrb}

\paragraph{Limits.} At $\lam = 0$, \cref{eq:pmix} is the Mie--Gr\"uneisen
pressure at $(\rho, e)$, and \bw\ evaluates it there with the inert
routine itself: a cell that has not started to burn is the inert material
bit for bit, tension and spall included. At $\lam = 1$ it is the ideal gas,
$P = (\gamma - 1)\rho(e + Q)$, and the solid is not evaluated at all.

\paragraph{Why burning can raise the pressure.} At equal density and
temperature a compressed solid is stiffer than the gas, but each unit of
$\lam$ adds $e_c + Q$ to $\thv$ (\cref{eq:theta}). Early in the reaction
the heat usually wins, so the pressure behind a shock can rise as the cell
burns, as it does in the reaction zone of a detonation.

\paragraph{A check with a closed form.} A cell at rest that burns
completely at constant volume ($\rho = \rhoz$, $e = 0$) ends at
$P = (\gamma - 1)\rhoz Q$ and $\thv = Q$: 12.499968~GPa for the class
problem's Material~X ($\rhoz = 1.60$~g/cm$^3$, $\gamma = 2.60$,
$Q = 4.8828$~kJ/g).
\trace{modules/mix_mod.f95}{mix_state}{test_hvrb}

\section{Tension: the solid dropped, then held}
\label{sec:tension}

\paragraph{Dropped.} A cell that has started to burn holds no tension: the
gas fills the cracks of a reacting solid. So where the solid would pull
($\Ps < 0$), it is left out of the pressure,
\begin{equation}
  P = (1 - \lam)\max(\Ps, 0) + \lam\,\Pg, \qquad 0 < \lam < 1.
  \label{eq:dropped}
\end{equation}
Only the solid's pressure is dropped; $\thv$, $T$ and the energy are as
\cref{eq:theta} writes them. A cell at $\lam = 0$ is not affected (it is the
inert material, which may spall).
\trace{modules/mix_mod.f95}{mix_state, mix_pulls}{test_hvrb}

\paragraph{Held where it first pulled.} A dropped solid does no work, yet as
the cell stretched its $e_c(\rho)$ went on changing, by
$\dd e_c = \Pc\,\dd\rho/\rho^2$ (\cref{eq:solid-theta}): energy moved into or
out of $\thv$ by no force. In runs this drove $\thv$, and $T$, negative
(to $-4300$~K, Material~X struck at 1~km/s with method AV). So \bw\ stops
the solid where it first pulls: it cavitates, and the gas fills the rest.
Each burning cell keeps a hold density $\rho_h$ and the isentrope's energy
there, $e_{S,h} = e_S(\rho_h)$, set once a step after $\lam$ and $e_S$ are
updated:
\begin{itemize}
\item while the cell is stretched past its hold ($\rho < \rho_h$): kept;
\item otherwise, if the solid pulls at the cell's own density: $\rho_h =
  \rho$, $e_{S,h} = e_S(\rho)$;
\item otherwise: not held ($\rho_h = 0$);
\end{itemize}
and never at $\lam = 0$ or $1$. While held, every term of the solid is
evaluated at $(\rho_h, e_{S,h})$ and the gas at the cell's own $\rho$:
$e_c$ is fixed at $e_{S,h} - \Cv T_0\exp(\Gz\eta_h)$, and
\begin{equation}
  P = s\,(1 - \lam)\,\Ps(\rho_h,\, e_c + \thv) + \lam\,(\gamma - 1)\rho\,\thv,
  \label{eq:held}
\end{equation}
with $s = 1$ if the held solid pushes and $s = 0$ if it would still pull.
At $\rho = \rho_h$ the held and unheld closures agree, so $P$ is
continuous; a cell that is never held gets the same bits as without the
rule. Held and dropped, an expansion at fixed $\lam$ cools the cell as
$\thv \propto \rho^{\lam(\gamma - 1)}$: never below zero, nor below the gas
alone expanded the same way.
\trace{modules/mix_mod.f95}{mix_state, mix_pulls}{test_hvrb}
\trace{modules/burn_mod.f95}{burn_hold}{run_hist_tests.sh}

\section{The sound speed}
\label{sec:sound-speed}

The sound speed at fixed $\lam$, of the pressure actually used, is
$c^2 = (\partial P/\partial\rho)_{e,\lam} + (P/\rho^2)(\partial P/\partial
e)_{\rho,\lam}$. It sets the time step and method AV's viscosity, and method
RS's impedance; it does not change the converged solution. For the closure
of \cref{eq:pmix} it is exactly
\begin{equation}
  c^2 = (1 - \lam)\,c_{\mathrm{sol}}^2 + \lam\,c_{\mathrm{gas}}^2
        - \lam(1 - \lam)\Bigl(\frac{\rhoz\Gz}{\rho} - \gamma + 1\Bigr)^{2}\thv,
  \label{eq:c-exact}
\end{equation}
with $c_{\mathrm{sol}}$ the inert Mie--Gr\"uneisen sound speed at $(\rho,
e_c + \thv)$: the mass-weighted mean of the phases' squares, less a term
from the spread of their two Gr\"uneisen coefficients ($\rhoz\Gz/\rho$ for
the solid, $\gamma - 1$ for the gas). The mass-weighted mean alone is about
1~\% high for Material~X; the equal-pressure (Wood) rule would be about
15~\% low.
\trace{modules/mix_mod.f95}{mix_state}{test_hvrb}

With the solid dropped (\cref{eq:dropped}), and held (\cref{eq:held}),
\begin{align}
  c^2_{\mathrm{dropped}} &= \lam(\gamma - 1)\Bigl[\thv\,\bigl(1 + \lam(\gamma - 1)\bigr)
     - (1 - \lam)\frac{\Pc}{\rho}\Bigr],
  \label{eq:c-dropped}\\
  c^2_{\mathrm{held}} &= \lam(\gamma - 1)\,\thv
     + \frac{P}{\rho^2}\bigl[s\,(1 - \lam)\rhoz\Gz + \lam(\gamma - 1)\rho\bigr].
  \label{eq:c-held}
\end{align}
Held and dropped ($s = 0$), \cref{eq:c-held} is $\lam(\gamma-1)\thv\,(1 +
\lam(\gamma - 1)) > 0$ wherever $\thv > 0$.

\paragraph{Guards.} If $c^2 \le 0$, or, for a cell that is not held,
$\thv \le 0$ (where the gas has no pressure and \cref{eq:c-exact} does not
hold), \bw\ falls back, in order, to the mass-weighted mean $(1 -
\lam)c_{\mathrm{sol}}^2 + \lam c_{\mathrm{gas}}^2$ (with $c_{\mathrm{sol}}$
at $\max(\Ps, 0)$, and at $\rho_h$ if the solid is held), then to the gas's
alone; if neither is positive the cell has no sound speed, and the run
stops as it does for an inert cell with none. Which rule gave $c$ is
returned with it.
\trace{modules/mix_mod.f95}{mix_state}{test_hvrb}

\section{The reaction source for method RS}
\label{sec:rs-source}

Method RS at second order moves each cell's pressure half a step inside the
step, in its predictor, with the acoustic equation $\dd P/\dd t =
-Z^2\,\partial u/\partial h$ ($Z = \rho c$, $h$ the mass coordinate). With
reaction the pressure equation gains a source, Fickett and Davis's
thermicity $\sigma$ times the rate $r = \dd\lam/\dd t$
\cite{fickett1979,kittell2020pep}:
\begin{equation}
  \frac{\dd P}{\dd t} = -Z^2\frac{\partial u}{\partial h} + \rho c^2\sigma\, r,
  \qquad
  \rho c^2\sigma = \Bigl(\frac{\partial P}{\partial\lam}\Bigr)_{\rho, e}.
  \label{eq:rs-source}
\end{equation}
For the closure, with the solid in and dropped,
\begin{align}
  \Bigl(\frac{\partial P}{\partial\lam}\Bigr)_{\rho,e}
  &= -\Ps + \Pg + \bigl[(1 - \lam)\rhoz\Gz + \lam(\gamma - 1)\rho\bigr](e_c + Q),
  \label{eq:dpdlam}\\
  \Bigl(\frac{\partial P}{\partial\lam}\Bigr)_{\rho,e,\ \mathrm{dropped}}
  &= \Pg + \lam(\gamma - 1)\rho\,(e_c + Q),
  \label{eq:dpdlam-dropped}
\end{align}
with $\gamma - 1$ taken as 0 where $\thv \le 0$ (no gas pressure there). At
$\lam = 1$ the dropped form is used (the solid is not evaluated, and nothing
multiplies it then); at $\lam = 0$ too, if the solid is in tension, since
the cell's first step of burning drops it (the pressure itself jumps there,
from $\Ps < 0$ to $P \ge 0$).
\trace{modules/mix_mod.f95}{mix_state}{test_hvrb}

Two choices make the source consistent with the rest of the step. It is held
constant over the step, from the cell's state at the step's start, as $Z$ is.
And its rate is the step's own, $r = (\lam^{n+1} - \lam^{n})/\Delta t$ from
the model's update, not the instantaneous rate at $\lam^n$ (HVRB's is
infinite at $\lam = 0$ for some exponents, and zero there for others; the
step's mean is always finite and is the $\lam$ the closure then uses). The
source shapes only the face pressures the Riemann problem is given; the
conservative update, and so energy conservation, are untouched. Order 1 has
no predictor and no source.
\trace{modules/burn_mod.f95}{burn_rate_rs}{compare_hvrb.py}
\trace{modules/rs_mod.f95}{rs_face_states}{compare_hvrb.py}

\section{Every cell's model, constants and history}
\label{sec:burn-table}

Every burn model's deck keys, per-cell constants and history variables are
described by one table (module \code{burn_table}); a new model is a row in
it, a module of its own for its rate, and a case in one routine. For cell
$k$, module \code{burn} keeps
\begin{itemize}
\item \code{bmod(k)}, the cell's model (0 if it does not burn);
\item \code{bpar(1:8, k)}, the constants of its rate law, in the slots its
  row names (8 is the most any of the six models needs);
\item \code{bst(1:6, k)}, its history variables (HVRB's $\phi$; the lead
  shock's five numbers);
\end{itemize}
and, for every model, $\lam$, the products' $\gamma$ and $Q$, and the hold
$(\rho_h, e_{S,h})$. Each cell's numbers lie together in memory, and every
cell writes only its own, so a run gives the same bits on any number of
threads. A deck without a burn allocates none of this and runs exactly as
before.
\trace{modules/burn_mod.f95}{alloc_burn, fill_constants}{test_btab, run_deck_tests.sh}

\paragraph{Keys.} Burn keys sit flat in the Material block, under three
rules: one key, one meaning, across all models (\key{PR}, the reference
pressure, serves every model that has one); each model accepts only its own
keys, and any other burn key is a deck error; and no burn key shares a name
with another Material key. \Cref{app:deck-keys} lists them.

\paragraph{The constants stay constant.} \code{bmod} and \code{bpar} are
written once, at setup, from the cell's layer. They are protected (no other
module can assign to them: it does not compile), and the models receive them
as \code{intent(in)} in pure routines. At the run's end \code{burn_audit}
fills them again from the layers and compares them bit by bit; a difference
stops the run with an error, since it could only be a fault in \bw. (Every
burning deck in \code{make test} runs this audit; that it does stop a run was
shown once, by changing a constant in a scratch build.) Graded
constants (planned) will vary from cell to cell, never in time, and change
only how \code{bpar} is filled.
\trace{modules/burn_mod.f95}{fill_constants, burn_audit}{run_deck_tests.sh}

\section{One cycle}
\label{sec:one-cycle}

Each model's rate uses the cell's state at the start of the step ($P^n$,
$T^n$, its history), and $\lam^{n+1}$ then feeds the closure. Where a
model's reaction topology has a closed form, its step is exact in $\lam$ for
that rate (\cref{ch:burn-models}); the only error is the rate's, first order
in $\Delta t$.

\paragraph{Method AV} (one cycle; \code{maincode}, section~7):
\begin{enumerate}[label=(\alph*)]
\item momentum, positions, volumes and energies, as for any cell;
\item the burn (\code{burn_step}): $e_S$ moved to the new volume, then each
  cell's model steps its history and $\lam$ from the start-of-step state
  (\code{cell_burn}), then the hold (\code{burn_hold},
  \cref{sec:tension});
\item the pressures: a cell with $\lam > 0$ from the closure at
  $(\rho, e, \lam)^{n+1}$, a cell at $\lam = 0$ as an inert cell;
\item the temperatures, $T = \thv/\Cv$ for a burning cell;
\item for the models that need it, the lead-shock detector, from each
  cell's state at the step's end (\code{burn_detect},
  \cref{sec:lead-shock}).
\end{enumerate}
Steps (d) and (e) close every cycle of method RS too.
\trace{modules/burn_mod.f95}{burn_step, cell_burn, burn_detect}{run_hist_tests.sh}

\paragraph{Method RS} (\code{maincode}, \code{rs_region}) needs $\lam^{n+1}$
before the predictor runs, so the burn comes first: the rate and the
reaction source from the state at the step's start
(\code{burn_rate_rs}, \cref{sec:rs-source}); then the predictor and the
Riemann problems; then each cell's update, $e_S$ to the new volume, the
hold, and the closure.

\paragraph{The gate.} A model that needs the lead shock
(\cref{sec:lead-shock}) does not burn until it is marked. The detector runs
at the end of a step, after the pressures, and the model reads the mark at
the start of the next, so a cell marked during step $n$ first burns in step
$n + 1$.

\paragraph{What a run reports.} The tracer file gets each tracer's $\lam$
(column \key{XRN}), the history file an \key{XRN} variable, and the run log
the first cell to react and the cells burnt out: those with $1 - \lam <
10^{-12}$, the same for every model (from 1.6.1; before, $\lam = 1$ exactly,
which ThermalCascade's step, only approaching 1, reached by rounding alone).

\section{Verification}
\label{sec:burning-cell-verification}

\begin{itemize}
\item \code{test_hvrb} checks the closure against 50-digit values
  (\code{tests/hvrb_reference.py}): $P$, $c$, $\thv$ and $\partial
  P/\partial\lam$ at states across the range, all by the exact sound speed;
  $\lam = 0$ is the inert routine bit for bit; the constant-volume burn ends
  at 12.499968~GPa with $\thv = Q$; the guards; and the hold (continuous at
  $\rho_h$, $c^2$ and $\partial P/\partial\lam$ equal to differences of $P$,
  and a held, dropped cell stretched to $\rho/100$ keeping $\thv > 0$ and
  following $\rho^{\lam(\gamma-1)}$).
\item The inert limit: a burn that never starts gives the inert run's
  heartbeat, tracer rows and history frames bit for bit, $\lam$'s own column
  aside (\code{tests/run_hist_tests.sh}, with the decks
  \code{ok_hvrb_inert} and \code{ok_tcas_inert}, methods AV and RS).
\item Energy: while Material~X burns, total energy is conserved to
  round-off by both methods (measured with the gold-standard comparisons,
  \code{tests/verification/compare_hvrb.py}, since 1.3.2;
  \code{doc/design/materials.md}). The script itself checks the
  reference's tolerance, 0.01~\%.
\item The closed forms of this chapter are derived from the definitions by
  \code{doc/manual/checks/ch6_closure.py} (sympy).
\end{itemize}
